\title{Regenerative Clamp Converter}
\subtitle{MS thesis · Center for Distributed Energy, Georgia Tech · Advisor: Dr.\ Deepak Divan}
\period{2025}
\date{2025-12}
\status{MS Thesis}
\order{1}
\tags{converters, gt}

\attach{MS thesis}{/thesis/Design_and_Analysis_of_a_Regenerative_Clamp_Converter_for_Energy_Capture_and_Recovery_in_High_Power_Converter_Topologies.pdf}

\begin{document}

Design and analysis of regenerative clamp converters that recover
leakage-inductor energy in high-power isolated converter topologies, reducing
device stress and improving efficiency. The clamp architecture was then extended
into a multiport interface, so the auxiliary port can integrate PV or another
distributed energy resource while still doing its job as a clamp.

\section{The problem}

In high-power isolated topologies, energy stored in the transformer's leakage
inductance has to go somewhere at every switching transition. A dissipative
clamp turns it into heat and voltage stress on the devices. Recovering that
energy instead of burning it buys back efficiency and lets the switches be
sized against a lower peak voltage.

\section{What the work covers}

\begin{itemize}
  \item Analysis and modelling of the regenerative clamp, including the
        conditions under which leakage energy is returned to the source rather
        than dissipated.
  \item Reduction of device voltage stress and the resulting effect on
        converter efficiency in high-power isolated topologies.
  \item Extension of the clamp into a multiport interface: the auxiliary port
        carries a PV or DER input while the clamping function is preserved.
  \item Simulation and hardware validation of the proposed architecture.
\end{itemize}

\end{document}
